% !TEX program = lualatex
\documentclass[12pt,a4paper]{article}
\usepackage{luatexja}
\usepackage{amsmath,amssymb,amsthm}
\usepackage{geometry}
\geometry{margin=1in}

%-------------------------------------------------
% マクロ
%-------------------------------------------------
\newcommand{\mweq}{\mathrel{\overset{\mathrm{mw}}{=}}}
\newcommand{\dfix}{D_{\mathrm{fix}0}}

%-------------------------------------------------
% 定理環境
%-------------------------------------------------
\newtheorem{definition}{定義}
\newtheorem{axiom}{公理}
\newtheorem{theorem}{定理}
\newtheorem{scholium}{補足}

\begin{document}

\title{空数学における二分木収束の存在証明}
\author{Masaomi Chiba}
\date{2026-04-20}
\maketitle

\begin{abstract}
本稿では、二分木分割列の収束を、構造的安定点 $\dfix$ に対する
中道等価 $\mweq$ によって記述する。
構造的歪み汎関数と中道等価の最小公理系を導入し、
二分木収束定理を示す。
\end{abstract}

\section{定義}

\begin{definition}[構造的歪み]
分割列 $(z_n)$ に対し、
\[
E_{n+1}=\frac12 E_n
\]
で定義する。
\end{definition}

\begin{definition}[歪み汎関数]
構造的安定点 $\dfix$ に対し
\[
D(z_n,\dfix):=E_n
\]
と定義する。
\end{definition}

\section{公理}

\begin{axiom}[中道等価導入]
もし
\[
D(x_n,\dfix)\to 0
\]
ならば
\[
x_n \mweq \dfix.
\]
\end{axiom}

\begin{axiom}[中道等価の反射律]
任意の $x$ に対し
\[
x \mweq x.
\]
\end{axiom}

\begin{axiom}[中道等価の対称律]
\[
x \mweq y
\quad\Rightarrow\quad
y \mweq x.
\]
\end{axiom}

\begin{axiom}[中道等価の推移律]
\[
x \mweq y
\;\land\;
y \mweq z
\quad\Rightarrow\quad
x \mweq z.
\]
\end{axiom}

\section{定理}

\begin{theorem}[二分木収束定理]
任意の二分木分割列 $(z_n)$ が

\[
E_{n+1}\le cE_n
\qquad (0<c<1)
\]

を満たすならば

\[
\lim_{n\to\infty}
z_n
\mweq
\dfix.
\]

\end{theorem}

\begin{proof}

定義より

\[
E_n=2^{-n}E_0.
\]

したがって

\[
\lim_{n\to\infty}E_n=0.
\]

歪み汎関数の定義より

\[
D(z_n,\dfix)=E_n
\]

なので

\[
D(z_n,\dfix)\to 0.
\]

ゆえに中道等価導入公理より

\[
\lim_{n\to\infty}
z_n
\mweq
\dfix.
\]

が成立する。

\end{proof}

\section{補足}

\begin{scholium}[Banach型類似]
本定理は
\[
E_{n+1}\le cE_n
\]
という縮小条件を通じて
\textit{Banach型}の収束構造を持つ。
ただし到達点ではなく
中道等価による安定化として解釈される。
\end{scholium}

\begin{scholium}[K\H{o}nig補題との関係]
本定理は
有限分岐無限木の存在論を
収束論的に再解釈する。
\end{scholium}

\end{document}
